\clearpage
\renewcommand{\thepage}{S\arabic{page}}
\setcounter{page}{1}
% Keep the supplementary title and wide task overview on the opening page.
\twocolumn[{%
    \centering
    {\Large\textbf{\thetitle}\par}
    \vspace{0.5em}
    {\Large Supplementary Material\par}
    \vspace{1em}
    \includegraphics[width=\textwidth]{assets/Robotwin.pdf}\par
    \captionsetup{type=figure,font=small,skip=4pt}
    \caption{Overview of the ten RoboTwin evaluation tasks.
    Task names and instructions are listed in
    Table~\ref{tab:robotwin-task-details}.}
    \label{fig:supp-simulation-task-overview}
    \vspace{1em}
}]

\section{Simulation Task Overview}
\label{app:robotwin-tasks}

We evaluate on ten tasks from RoboTwin 2.0~\cite{chen2025robotwin2},
illustrated in Fig.~\ref{fig:supp-simulation-task-overview}.
Table~\ref{tab:robotwin-task-details} lists their names, representative
instructions, and average episode lengths.

\begin{table}[!ht]
    \centering
    \caption{Overview of all ten RoboTwin evaluation tasks. Avg. steps
    denotes average recorded demonstration length, not the rollout limit.
    Instructions summarize the task goals.}
    \label{tab:robotwin-task-details}
    \footnotesize
    \setlength{\tabcolsep}{2.5pt}
    \renewcommand{\arraystretch}{1.12}
    \begin{tabular}{@{}
        >{\raggedright\arraybackslash}p{0.31\columnwidth}
        >{\raggedleft\arraybackslash}p{0.085\columnwidth}
        >{\raggedright\arraybackslash}p{\dimexpr0.605\columnwidth-4\tabcolsep\relax}
        @{}}
        \toprule
        Task & \shortstack{Avg.\\steps} & Instruction \\
        \midrule
        Shake Bottle Horizontally & 276
        & Pick up the bottle and shake it horizontally. \\
        Place Container Plate & 156
        & Place the container on the plate. \\
        Stack Blocks Two & 316
        & Move the red block to the center and stack the green block on it. \\
        Dump Bin Bigbin & 265
        & Pour the balls from the small bin into the large bin. \\
        Stack Bowls Three & 476
        & Stack the three bowls on top of one another. \\
        Pick Dual Bottles & 127
        & Pick up one bottle with each arm simultaneously. \\
        Handover Mic & 223
        & Pick up the microphone and hand it to the other arm. \\
        Place Bread Skillet & 162
        & Hold the skillet with one arm and place the bread inside with the other. \\
        Place Dual Shoes & 228
        & Place both shoes in the shoebox with their tips pointing left. \\
        Put Bottles Dustbin & 637
        & Put all three bottles into the dustbin on the left. \\
        \bottomrule
    \end{tabular}
\end{table}

% Task names/instructions were checked against official RoboTwin environments:
% https://github.com/RoboTwin-Platform/RoboTwin
% Documentation snapshot: ea8b21121ebb3cd201ff5b3fe361944ac94eda3f.
% The snapshot above is not asserted to be the evaluation checkout revision.
% Average lengths are retained from the author's original ten-task summary;
% no new episode statistics have been computed in this edit.
% Environment identifiers, in table order:
% shake_bottle_horizontally, place_container_plate, stack_blocks_two,
% dump_bin_bigbin, stack_bowls_three, pick_dual_bottles, handover_mic,
% place_bread_skillet, place_dual_shoes, put_bottles_dustbin.

\subsection{Evaluation Protocol}
\label{app:robotwin-eval-protocol}

Each method is evaluated with 100 rollouts per task using different random
seeds. We use RoboTwin's task-specific success checks and report both
per-task success and the unweighted mean across the ten tasks.
Policies predict 32-step action chunks, execute the first 16 steps, and
then replan from a new observation. All ten evaluation tasks belong to
the 50-task collection used for training.

\section{Implementation Details}
\label{app:baselines}

This section complements Sec.~\ref{sec:benchmark-and-baselines} with the
simulation training details of BiCAT and the baselines.

\subsection{Training Data and Inputs}
\label{app:baseline-data}

The 5\% validation split is drawn at the episode level with a fixed seed,
giving 2,375 training and 125 validation episodes, and is shared by
tokenizer and policy training. Action chunks of 32 steps are sampled within
episode boundaries; each step has 14 dimensions, six joint values and one
gripper value per arm. The three RGB views come from the head and the two
wrist cameras, and the robot state is encoded as a text prefix with 256
bins per dimension. All policies use the
\texttt{paligemma2-3b-mix-224} checkpoint.

\subsection{Optimization Settings}
\label{app:baseline-protocol}

Table~\ref{tab:supp-training-settings} lists the optimization settings
shared by all methods. Policies minimize teacher-forced action-token
cross-entropy and use greedy decoding during evaluation.

\begin{table}[!ht]
    \centering
    \caption{Shared optimization settings for the simulation comparison.
    Tokenizer pretraining applies to QueST, both OAT variants,
    $\mathcal{M}^{2}$Tok, NAC, ActionCodec, and BiCAT. The policy learning rate applies equally to
    the language model, vision encoder, and modality bridge.}
    \label{tab:supp-training-settings}
    \footnotesize
    \setlength{\tabcolsep}{3pt}
    \renewcommand{\arraystretch}{1.15}
    \begin{tabular}{@{}
        >{\raggedright\arraybackslash}p{0.40\columnwidth}
        >{\centering\arraybackslash}p{0.30\columnwidth}
        >{\centering\arraybackslash}p{\dimexpr0.30\columnwidth-4\tabcolsep\relax}
        @{}}
        \toprule
        Setting & \shortstack{Tokenizer\\pretraining}
        & \shortstack{Policy\\training} \\
        \midrule
        Optimizer steps & 300k & 30k \\
        Batch size & 512 & 64 \\
        Optimizer & AdamW & AdamW \\
        Learning rate & $5\times10^{-5}$ & $10^{-4}$ \\
        Adam $\boldsymbol{\beta}$ & $(0.9,0.95)$ & $(0.95,0.999)$ \\
        Adam $\epsilon$ & $10^{-8}$ & $10^{-6}$ \\
        Weight decay & 0 & $10^{-6}$ \\
        Warmup steps & 200 & 500 \\
        Schedule after warmup & Constant & Cosine \\
        Final learning rate & $5\times10^{-5}$ & $10^{-5}$ \\
        Precision & BF16 & BF16 \\
        Gradient norm limit & 1.0 & 1.0 \\
        Validation interval & 1k steps & 1k steps \\
        \bottomrule
    \end{tabular}
\end{table}

\subsection{Tokenizer Configurations}
\label{app:tokenizer-configurations}

Table~\ref{tab:efficiency} in the main paper lists the token count,
vocabulary size, and prediction schedule of each tokenizer. All learned
tokenizers use a hidden width of 256 unless noted and are trained on action
reconstruction.
\begin{itemize}
    \item \textbf{Bin}~\cite{kim2024openvla}: each normalized action value is
    mapped to one of 256 uniform bins.
    \item \textbf{FAST}~\cite{pertsch2025fast}: a discrete cosine transform,
    coefficient quantization, and byte-pair encoding, fitted on the RoboTwin
    training split rather than taken from the released universal tokenizer;
    it yields 26.8 tokens per chunk on average (range 9--80).
    \item \textbf{QueST}~\cite{mete2024quest}: causal temporal convolutions
    followed by a six-layer causal Transformer encoder, with FSQ levels
    $[8,8,6,5]$.
    \item \textbf{OAT}~\cite{liu2026oat_rss}: learned register tokens trained
    with nested prefix dropout over $\{1,2,4,8,16\}$ registers, with FSQ
    levels $[8,8,6,5]$; the two variants differ only in their register and
    prediction block schedules.
    \item $\boldsymbol{\mathcal{M}^{2}}$\textbf{Tok}~\cite{xu2026m2tok}: the
    released implementation with width 120. A shared convolutional VQ-VAE
    encodes each arm independently into two latent frames, each quantized by
    four cosine-similarity codebooks of 512 entries, and reconstructs each
    arm from its own tokens.
    \item \textbf{NAC}~\cite{jawaid2026nac}: a SEANet encoder over the
    time-major flattened chunk and two-level residual FSQ with levels
    $[8,8,4,4]$, a coarse level of 8 tokens and a fine level of 16; its
    Vocos decoder is trained with reconstruction, adversarial, and
    feature-matching losses.
    \item \textbf{ActionCodec}~\cite{dong2026actioncodec}: a six-layer
    Perceiver-style encoder and decoder, with 16 latent vectors quantized by
    a 2048-entry VQ codebook updated by exponential moving average.
    \item \textbf{BiCAT}: the architecture of Sec.~\ref{sec:method}, with one
    shared causal self-attention layer and one symmetric cross-arm attention
    layer in the encoder, FSQ levels $[8,8,6,5]$, and a six-layer joint
    decoder.
\end{itemize}

\section{Per-Task Ablation Results}
\label{app:ablation-per-task}

Table~\ref{tab:supp-ablation-per-task} reports per-task success for the
variants summarized in Table~\ref{tab:ablation}.


\section{Shared Arm Representations}
\label{app:arm-latents}

Figure~\ref{fig:arm-latents} visualizes the pre-quantization latents of
BiCAT, \emph{Raw coordinates}, and \emph{Arm-specific modules}. We encode
10,000 action chunks from the ten evaluation tasks, average each arm's eight
pre-quantization latent vectors, and embed the left- and right-arm averages
of each model jointly with t-SNE. With all modules shared, the latents of
the two arms overlap. With \emph{Arm-specific modules}, the two arms occupy
largely separate regions, so the shared FSQ grid assigns the same token IDs
to latents that are distributed differently for the two arms.

We quantify the overlap with the maximum mean discrepancy (MMD) between the
left- and right-arm averages in the original 4-D latent space, using an RBF
kernel whose bandwidth is the median pairwise squared distance. The MMD is
0.088 for BiCAT, 0.193 for \emph{Raw coordinates}, and 0.286 for
\emph{Arm-specific modules}, and the ordering is unchanged after pooled
per-dimension standardization. Under \emph{Raw coordinates}, the two arms
remain visually mixed in the t-SNE embedding, yet their MMD is more than
twice that of BiCAT, so canonicalization aligns the arms beyond what the
shared encoder alone achieves. These distances describe arm-marginal
distributions: they are consistent with, but do not directly measure, a
token ID denoting the same motion for both arms, and they do not by
themselves measure coordination or task success.


\section{Per-Task FSQ Ablation}
\label{app:fsq-per-task}

Table~\ref{tab:supp-fsq-per-task} reports per-task success for the FSQ
grids summarized in Table~\ref{tab:FSQ}. All grids use the BiCAT
architecture and the common training and evaluation protocol.


\section{Real-World Experiment Details}
\label{app:real-world}


Table~\ref{tab:real-world-task-details} lists the six real-world tasks.
All demonstrations are collected by teleoperation and recorded at 30\,Hz,
which is also the control frequency at deployment.

\subsection{Training Settings}
\label{app:real-world-training}

\begin{table}[!ht]
    \centering
    \caption{Real-world optimizer steps per embodiment. All tokenizers on
    the same embodiment use the same numbers of steps.}
    \label{tab:real-world-training}
    \footnotesize
    \setlength{\tabcolsep}{6pt}
    \renewcommand{\arraystretch}{1.12}
    \begin{tabular}{@{}lcc@{}}
        \toprule
        Embodiment & Tokenizer steps & Policy steps \\
        \midrule
        Bimanual SO-101 & 50k  & 4k \\
        Unitree G1     & 80k  & 9k \\
        Bimanual YAM   & 300k & 20k \\
        \bottomrule
    \end{tabular}
\end{table}

Each embodiment is trained separately. For each embodiment, one action
tokenizer is trained on the demonstrations of all its tasks in
Table~\ref{tab:real-world-task-details} without a held-out validation
split and then frozen, and one VLA policy is trained on its
tokens jointly for all of these tasks; on Unitree G1, for example, a single
tokenizer and a single policy cover both Bottle and Bowl.
Table~\ref{tab:real-world-training} lists the numbers of optimizer steps,
which scale with the amount of data on each embodiment. BiCAT and
$\mathrm{OAT}^{\mathrm{sing}}$ use the same tokenizer architectures as in
simulation, with 16 tokens per chunk and FSQ levels $[8,8,6,5]$, and FAST is
fitted separately on the demonstrations of each embodiment.

\paragraph{Policy.}
The real-world policies use Molmo2-ER~\cite{fang2026molmoact2} as the
vision-language backbone, which is fully fine-tuned together with the
vision encoder and the modality bridge. The policy observes the language
instruction, the joint state of both arms, and three RGB views from the
current time step: a top camera and the two wrist cameras. On Bimanual YAM
and Unitree G1 the top view has a resolution of $360\times640$ and the
wrist views $480\times640$; on Bimanual SO-101 all three views are
$480\times640$. The policy predicts 32-step action chunks as in
simulation, but on the real robots it executes all 32 steps of each chunk
before replanning, rather than the first 16.

\paragraph{Optimization.}
Tokenizer pretraining uses the settings of
Table~\ref{tab:supp-training-settings} except for the number of steps
(Table~\ref{tab:real-world-training}). Policy training uses the same
peak learning rate and cosine decay as in simulation, with a global batch
size of 128, BF16 mixed precision, and gradient-norm clipping at 1.0.
Within an embodiment, the training configuration is identical across
methods: BiCAT and $\mathrm{OAT}^{\mathrm{sing}}$ are trained for the same
number of tokenizer steps, FAST is fitted on the same demonstrations, and
all three policies share the data, observations, number of policy steps,
learning rate, warmup, and batch size. They differ only in the tokenizer
and its prediction schedule.

\subsection{Evaluation Protocol}
\label{app:real-world-eval}

Each method is evaluated for 25 trials per task. A trial counts as a
success only if the robot completes every stage of the task without error;
on Block, for example, the block must be placed in the box and the box
handed over. Partial completion counts as a failure. During deployment, the
robot waits while the policy predicts the next action chunk and executes
the chunk once inference has finished, so differences in inference latency
lengthen the episode but do not change the observation on which each chunk
is predicted. The initial object placements are identical across methods
for each trial, and each trial is limited to 120\,s, after which it counts
as a failure. Inference latency is measured from receiving an observation to outputting
an action chunk, over all inference calls during the YAM evaluation, and is
reported as mean $\pm$ standard deviation.

% Full-width appendix floats are collected here so that they do not interrupt the two-column text.
\begin{table*}[!ht]
    \centering
    \caption{Per-task success rates of the BiCAT ablations on the ten
    RoboTwin 2.0 tasks, evaluated over 100 rollouts per task.}
    \label{tab:supp-ablation-per-task}
    \footnotesize
    \setlength{\tabcolsep}{3.5pt}
    \resizebox{\textwidth}{!}{%
    \begin{tabular}{@{}l*{10}{c}c@{}}
        \toprule
        Variant
        & \shortstack[c]{Shake bottle\\horizontally}
        & \shortstack[c]{Place container\\plate}
        & \shortstack[c]{Stack blocks\\two}
        & \shortstack[c]{Dump bin\\bigbin}
        & \shortstack[c]{Stack bowls\\three}
        & \shortstack[c]{Pick dual\\bottles}
        & \shortstack[c]{Hand over\\mic}
        & \shortstack[c]{Place bread\\skillet}
        & \shortstack[c]{Place dual\\shoes}
        & \shortstack[c]{Put bottles\\dustbin}
        & Avg. \\
        \midrule
        Raw coordinates & 0.93 & 0.88 & 0.66 & 0.85 & 0.64 & 0.22 & 0.95 & 0.26 & 0.24 & 0.24 & 0.587 \\
        Arm-specific modules & 0.94 & 0.87 & 0.51 & 0.85 & 0.75 & 0.21 & 0.94 & 0.24 & 0.05 & 0.33 & 0.569 \\
        \midrule
        Decoder-only coupling & 0.95 & 0.73 & 0.44 & 0.68 & 0.55 & 0.12 & 0.96 & 0.31 & 0.10 & 0.33 & 0.517 \\
        Encoder-only coupling & 0.96 & 0.76 & 0.47 & 0.62 & 0.65 & 0.31 & 0.92 & 0.42 & 0.29 & 0.20 & 0.560 \\
        Independent arms & 0.89 & 0.82 & 0.37 & 0.73 & 0.58 & 0.34 & 0.70 & 0.38 & 0.08 & 0.11 & 0.500 \\
        \midrule
        Single-block prediction & 0.90 & 0.89 & 0.61 & 0.68 & 0.64 & 0.17 & 0.62 & 0.21 & 0.06 & 0.33 & 0.511 \\
        Right-first prediction & 0.97 & 0.87 & 0.47 & 0.84 & 0.61 & 0.17 & 0.92 & 0.37 & 0.04 & 0.16 & 0.542 \\
        \midrule
        \oursrow\textbf{BiCAT} & 1.00 & 0.92 & 0.66 & 0.85 & 0.81 & 0.44 & 0.88 & 0.47 & 0.41 & 0.49 & 0.693 \\
        \bottomrule
    \end{tabular}%
    }
\end{table*}

\begin{figure*}[!ht]
    \centering
    \newcommand{\armpanel}[3]{%
        \begin{minipage}[t]{0.32\textwidth}\centering
        {\small #1}\\[2pt]
        \includegraphics[width=\linewidth,trim=#2 53.6 #3 27,clip]{assets/arm_mean_pre_fsq.pdf}
        \end{minipage}}
    \armpanel{BiCAT}{18}{519.7}\hfill
    \armpanel{\emph{Raw coordinates}}{268.4}{269.2}\hfill
    \armpanel{\emph{Arm-specific modules}}{519}{18.6}\\[3pt]
    {\small\textcolor[RGB]{208,88,88}{\large$\bullet$}~Left arm\qquad
    \textcolor[RGB]{88,144,208}{\large$\bullet$}~Right arm}
    \caption{t-SNE of pre-quantization arm latents on the ten RoboTwin 2.0
    tasks. Each point is one arm of one action chunk, averaged over its eight
    latent positions; each panel contains the same 10,000 chunks. Red and
    blue shades denote the left and right arms. Each panel is fitted
    separately, so coordinates and distances are not comparable across
    panels; the left--right MMD reported in the text is computed in the
    original latent space.}
    \label{fig:arm-latents}
\end{figure*}

\begin{table*}[!ht]
    \centering
    \caption{Per-task success rates of BiCAT with different FSQ levels on
    the ten RoboTwin 2.0 tasks, evaluated over 100 rollouts per task.
    $|\mathcal{V}|$ is the induced vocabulary size, and Avg.\ is the
    unweighted task average.}
    \label{tab:supp-fsq-per-task}
    \footnotesize
    \setlength{\tabcolsep}{3.5pt}
    \resizebox{\textwidth}{!}{%
    \begin{tabular}{@{}lc*{10}{c}c@{}}
        \toprule
        FSQ levels & $|\mathcal{V}|$
        & \shortstack[c]{Shake bottle\\horizontally}
        & \shortstack[c]{Place container\\plate}
        & \shortstack[c]{Stack blocks\\two}
        & \shortstack[c]{Dump bin\\bigbin}
        & \shortstack[c]{Stack bowls\\three}
        & \shortstack[c]{Pick dual\\bottles}
        & \shortstack[c]{Hand over\\mic}
        & \shortstack[c]{Place bread\\skillet}
        & \shortstack[c]{Place dual\\shoes}
        & \shortstack[c]{Put bottles\\dustbin}
        & Avg. \\
        \midrule
        $[8,6,5]$     & 240  & 0.87 & 0.86 & 0.47 & 0.69 & 0.51 & 0.12 & 0.82 & 0.12 & 0.04 & 0.20 & 0.470 \\
        $[8,8,8]$     & 512  & 0.99 & 0.84 & 0.65 & 0.84 & 0.70 & 0.22 & 0.97 & 0.35 & 0.21 & 0.40 & 0.617 \\
        $[8,5,5,5]$   & 1000 & 0.92 & 0.91 & 0.61 & 0.84 & 0.71 & 0.21 & 0.97 & 0.42 & 0.21 & 0.51 & 0.631 \\
        $[8,8,6,5]$   & 1920 & 1.00 & 0.92 & 0.66 & 0.85 & 0.81 & 0.44 & 0.88 & 0.47 & 0.41 & 0.49 & 0.693 \\
        $[7,5,5,5,5]$ & 4375 & 0.90 & 0.85 & 0.51 & 0.81 & 0.61 & 0.34 & 0.91 & 0.34 & 0.23 & 0.39 & 0.589 \\
        \bottomrule
    \end{tabular}%
    }
\end{table*}

\begin{table*}[!ht]
    \centering
    \caption{Real-world tasks. The task description summarizes each task and
    need not match the language instruction given to the policy. Action dim.\
    is the per-arm action dimension times the number of arms. Demos is the
    number of teleoperated demonstrations, and Avg.\ steps is the average
    number of recorded time steps per demonstration.}
    \label{tab:real-world-task-details}
    \footnotesize
    \setlength{\tabcolsep}{4pt}
    \renewcommand{\arraystretch}{1.12}
    \begin{tabular}{@{}llp{0.42\textwidth}cccc@{}}
        \toprule
        Embodiment & Task & Task description & Action dim. & Demos & Avg.\ steps & Rec.\ rate \\
        \midrule
        Bimanual SO-101 & Block  & Put the block into the box, then hand the box over. & $6\times2$ & 100 & 544 & 30\,Hz \\
        \midrule
        Unitree G1 & Bottle & Move the bottle to the left. & $8\times2$ & 100 & 441 & 30\,Hz \\
                   & Bowl   & Push the plate to the center, then place the bowl on it. & $8\times2$ & 100 & 538 & 30\,Hz \\
        \midrule
        Bimanual YAM & Pen    & Grasp the marker and its cap, then close the marker. & $7\times2$ & 105 & 447 & 30\,Hz \\
                     & Box    & Put all the metal parts in the box, then move the box forward. & $7\times2$ & 100 & 733 & 30\,Hz \\
                     & Bottle & Put all bottles in the dustbin. & $7\times2$ & 100 & 1236 & 30\,Hz \\
        \bottomrule
    \end{tabular}
    % Frames from the Hugging Face dataset cards (checked 2026-11-12):
    % SO-101 block 54,412; G1 bottles 44,123, plate 53,822;
    % YAM pen 46,924, box 73,252, bottles 123,558. Avg. steps = frames / demos.
    % Action dims: SO-101 5 joints + gripper; G1 7 joints + gripper; YAM 6 joints + gripper.
\end{table*}
